\chapter{Criterio de admisión metodológica}

\section{Admisión al nivel metodológico}
\begin{definition}
Se llama admisión metodológica al acto disciplinado mediante el cual una pieza preparada en el canon intermedio es aceptada para tratamiento en Tomo III.
\end{definition}

\begin{observation}
La admisión metodológica no convierte por sí sola a la pieza en demo consolidada ni en claim final. Solo declara que ha alcanzado un umbral suficiente para entrar en evaluación selectiva, curatorial o demostrativa.
\end{observation}

\section{Condición mínima de admisión}
\begin{definition}
Una pieza puede ser admitida metodológicamente cuando concurren simultáneamente:
\begin{enumerate}
  \item dependencia canónica visible;
  \item estabilización operativa suficiente en Tomo II;
  \item protocolización mínima legible;
  \item trazabilidad operativa suficiente;
  \item función reconocible dentro de un posible uso metodológico, demostrativo o selectivo.
\end{enumerate}
\end{definition}

\begin{observation}
La admisión no exige todavía validación exhaustiva ni clausura completa de resultados. Exige solo que la pieza ya no sea prematura para tratamiento metodológico.
\end{observation}

\section{Admisión y configuración}
\begin{definition}
Una configuración operativa es admisible en Tomo III cuando puede ser presentada como unidad tratable, con función, origen y recorrido suficientemente legibles para uso metodológico.
\end{definition}

\begin{observation}
Esta cláusula evita que el Tomo III reciba configuraciones cuya identidad siga siendo opaca o demasiado inestable.
\end{observation}

\section{Admisión y familia}
\begin{definition}
Una familia operativa es admisible en Tomo III cuando sus miembros pueden ser reunidos como conjunto disciplinado, con núcleo común reconocible y diferencias internas no destructivas para la finalidad metodológica.
\end{definition}

\begin{observation}
La admisión de una familia no exige uniformidad total. Exige cohesión suficiente para tratamiento selectivo o demostrativo.
\end{observation}

\section{Admisión y traducción}
\begin{definition}
Una traducción entre niveles es admisible en Tomo III cuando puede explicitar qué conserva, qué modifica y bajo qué condiciones sigue siendo metodológicamente utilizable.
\end{definition}

\begin{observation}
Esta cláusula impide admitir traducciones cuya utilidad dependa de saltos implícitos o pérdidas no declaradas.
\end{observation}

\section{Condición mínima de no admisión}
\begin{definition}
Se considera que una pieza no es metodológicamente admisible cuando falla de forma significativa en dependencia visible, estabilización operativa, protocolización mínima, trazabilidad suficiente o función reconocible.
\end{definition}

\begin{observation}
Esta condición protege al Tomo III frente a material prematuro, híbrido o insuficientemente consolidado.
\end{observation}

\section{Cláusula de alcance}
\begin{observation}
El criterio de admisión metodológica fija la frontera mínima de entrada a Tomo III. No constituye todavía una teoría completa de auditoría, selección final de claims o validación exhaustiva.
\end{observation}

